% --- Archon physics formalization source begin ---
% archon:physics
% archon:covers PhyXMiniProblems/problem_phyx_mini_0111.lean
% archon:source-report reports/phyx_mini/problem_phyx_mini_0111.source.json

\chapter{Physics problem phyx\_mini\_0111}
\label{ch:PhyXMiniProblems_problem_phyx_mini_0111}

\paragraph{Problem source.}
\#\# Physical scenario

Figure is a diagram of a Galilean telescope, or opera glass, with both the object and its final image at infinity. The image \$I\$ serves as a virtual object for the eyepiece. The final image is virtual and erect. The Galilean telescope has an angular magnification of -6.33.

\#\# Question

What focal length should the eyepiece have?

\#\# Answer choices

- A: A: -15.0cm

- B: B: -14.0cm

- C: C: -13.0cm

- D: D: -12.0cm

\#\# Figure caption (auxiliary; use the image as primary evidence)

The image is an optical diagram depicting the components and path of light through a lens system, likely a telescope or microscope. Here's a detailed breakdown:

1. **Components:**
   - **Objective Lens**: A convex lens on the left, labeled “Objective.”
   - **Eyepiece Lens**: A concave lens on the right, labeled “Eyepiece.”

2. **Light Paths:**
   - Purple lines represent the light rays entering from the left, passing through the objective lens, and exiting through the eyepiece lens.

3. **Focal Points and Distances:**
   - **\textbackslash{}(F\_1'\textbackslash{}) and \textbackslash{}(F\_2\textbackslash{})**: Focal points located along the horizontal optical axis near the eyepiece.
   - **\textbackslash{}(F\_2\textbackslash{})**: Also near the objective lens.
   - Distances are annotated as \textbackslash{}(f\_1\textbackslash{}) and \textbackslash{}(f\_2\textbackslash{}) along the optical axis, indicating focal lengths or distances between the focal points and lenses.

4. **Parallel Rays and Angles:**
   - Incoming rays are parallel before converging through the objective lens.
   - Rays diverge after passing through the eyepiece, with dashed lines indicating the paths if extended backward.

5. **Image Formation:**
   - A point labeled “I,” indicating the image position along the optical axis after light exits the eyepiece.

6. **Labels and Notations:**
   - Vectors are labeled \textbackslash{}( \textbackslash{}mathbf\{u\} \textbackslash{}) and \textbackslash{}( \textbackslash{}mathbf\{u'\} \textbackslash{}), showing the direction and magnitude of rays entering and exiting the objective lens.
   - The diagram includes arrows showing the direction of light travel.

The image is a typical diagram used to illustrate the function of compound lens systems in devices like telescopes and microscopes.

\#\# Dataset metadata

Optical Instruments; Physical Model Grounding Reasoning, Multi-Formula Reasoning

\paragraph{Recorded answer/context.}
A: A: -15.0cm

\paragraph{Figure/image path.}
/root/proposal\_for\_physic/science-mango/phyx\_mini\_run/phyx\_data/test\_image/111.png

\paragraph{Formalization target.}
create a compiling Lean file with sorry bodies at `PhyXMiniProblems/problem\_phyx\_mini\_0111.lean`. The Lean declarations must preserve the physical quantities, dimensions or dimensional roles, figure labels, governing-law hypotheses, and final relation expressed by this problem.
Use Mathlib/Physlib names found through LeanExplore where available. If a physics API is missing, introduce faithful local abstractions rather than scalar placeholder aliases.

\begin{theorem}[Physics formalization target]
\label{thm:physics:phyx_mini_0111:target}
\lean{PhyXMiniProblems.ProblemPhyXMini0111.problem_phyx_mini_0111}
\uses{def:physics:phyx-mini-0111:phyxminiproblems-problemphyxmini0111-lengthreadout, def:physics:phyx-mini-0111:phyxminiproblems-problemphyxmini0111-galileantelescopesetup, def:physics:phyx-mini-0111:phyxminiproblems-problemphyxmini0111-matchesgalileantelescopefigure, def:physics:phyx-mini-0111:phyxminiproblems-problemphyxmini0111-matchesangularmagnificationreadout, def:physics:phyx-mini-0111:phyxminiproblems-problemphyxmini0111-hasphysicalgalileanparameters, def:physics:phyx-mini-0111:phyxminiproblems-problemphyxmini0111-satisfiesafocalgalileanlaws}
The assigned autoformalize agent should translate this physics problem into Lean declarations in the covered file, with theorem and lemma proof bodies written as `by sorry`.
\end{theorem}
\begin{proof}
This is an autoformalization task, not a proof task. Produce statements that can later be proved by the physics prover without weakening the physical model.
\end{proof}

% --- Archon declaration topology begin ---
\section{Declaration topology}
\begin{definition}[Optical Length]
  \label{def:physics:phyx-mini-0111:phyxminiproblems-problemphyxmini0111-opticallength}
  \lean{PhyXMiniProblems.ProblemPhyXMini0111.OpticalLength}
  A signed physical quantity carrying the dimension of length
\end{definition}
\begin{proof}
  This is the declared model definition of Optical Length.
\end{proof}
\begin{definition}[length Readout]
  \label{def:physics:phyx-mini-0111:phyxminiproblems-problemphyxmini0111-lengthreadout}
  \lean{PhyXMiniProblems.ProblemPhyXMini0111.lengthReadout}
  \uses{def:physics:phyx-mini-0111:phyxminiproblems-problemphyxmini0111-opticallength}
  The scalar readout of a physical length in a chosen length unit
\end{definition}
\begin{proof}
  This is the declared model definition of length Readout.
\end{proof}
\begin{definition}[centimeters Value]
  \label{def:physics:phyx-mini-0111:phyxminiproblems-problemphyxmini0111-centimetersvalue}
  \lean{PhyXMiniProblems.ProblemPhyXMini0111.centimetersValue}
  \uses{def:physics:phyx-mini-0111:phyxminiproblems-problemphyxmini0111-opticallength, def:physics:phyx-mini-0111:phyxminiproblems-problemphyxmini0111-lengthreadout}
  The scalar readout of a physical length in centimeters
\end{definition}
\begin{proof}
  This is the declared model definition of centimeters Value.
\end{proof}
\begin{definition}[Lens Role]
  \label{def:physics:phyx-mini-0111:phyxminiproblems-problemphyxmini0111-lensrole}
  \lean{PhyXMiniProblems.ProblemPhyXMini0111.LensRole}
  Roles of the two lenses along the telescope's direction of propagation
\end{definition}
\begin{proof}
  This is the declared model definition of Lens Role.
\end{proof}
\begin{definition}[Thin Lens Kind]
  \label{def:physics:phyx-mini-0111:phyxminiproblems-problemphyxmini0111-thinlenskind}
  \lean{PhyXMiniProblems.ProblemPhyXMini0111.ThinLensKind}
  Paraxial lens kind, including the sign information of its optical power
\end{definition}
\begin{proof}
  This is the declared model definition of Thin Lens Kind.
\end{proof}
\begin{definition}[Thin Lens]
  \label{def:physics:phyx-mini-0111:phyxminiproblems-problemphyxmini0111-thinlens}
  \lean{PhyXMiniProblems.ProblemPhyXMini0111.ThinLens}
  \uses{def:physics:phyx-mini-0111:phyxminiproblems-problemphyxmini0111-opticallength, def:physics:phyx-mini-0111:phyxminiproblems-problemphyxmini0111-thinlenskind}
  A thin lens with a signed physical focal length
\end{definition}
\begin{proof}
  This is the declared model definition of Thin Lens.
\end{proof}
\begin{definition}[Conjugate Location]
  \label{def:physics:phyx-mini-0111:phyxminiproblems-problemphyxmini0111-conjugatelocation}
  \lean{PhyXMiniProblems.ProblemPhyXMini0111.ConjugateLocation}
  Whether an object or image conjugate is at a finite axial point or infinity
\end{definition}
\begin{proof}
  This is the declared model definition of Conjugate Location.
\end{proof}
\begin{definition}[Image Nature]
  \label{def:physics:phyx-mini-0111:phyxminiproblems-problemphyxmini0111-imagenature}
  \lean{PhyXMiniProblems.ProblemPhyXMini0111.ImageNature}
  Whether rays physically meet at an image or only appear to originate there
\end{definition}
\begin{proof}
  This is the declared model definition of Image Nature.
\end{proof}
\begin{definition}[Image Orientation]
  \label{def:physics:phyx-mini-0111:phyxminiproblems-problemphyxmini0111-imageorientation}
  \lean{PhyXMiniProblems.ProblemPhyXMini0111.ImageOrientation}
  Transverse image orientation relative to the observed object
\end{definition}
\begin{proof}
  This is the declared model definition of Image Orientation.
\end{proof}
\begin{definition}[Axial Figure Point]
  \label{def:physics:phyx-mini-0111:phyxminiproblems-problemphyxmini0111-axialfigurepoint}
  \lean{PhyXMiniProblems.ProblemPhyXMini0111.AxialFigurePoint}
  Axial points explicitly displayed in the Galilean-telescope figure
\end{definition}
\begin{proof}
  This is the declared model definition of Axial Figure Point.
\end{proof}
\begin{definition}[Ray Direction Label]
  \label{def:physics:phyx-mini-0111:phyxminiproblems-problemphyxmini0111-raydirectionlabel}
  \lean{PhyXMiniProblems.ProblemPhyXMini0111.RayDirectionLabel}
  Labels of the incident and emergent ray directions shown in the figure
\end{definition}
\begin{proof}
  This is the declared model definition of Ray Direction Label.
\end{proof}
\begin{definition}[Galilean Telescope Setup]
  \label{def:physics:phyx-mini-0111:phyxminiproblems-problemphyxmini0111-galileantelescopesetup}
  \lean{PhyXMiniProblems.ProblemPhyXMini0111.GalileanTelescopeSetup}
  \uses{def:physics:phyx-mini-0111:phyxminiproblems-problemphyxmini0111-opticallength, def:physics:phyx-mini-0111:phyxminiproblems-problemphyxmini0111-lensrole, def:physics:phyx-mini-0111:phyxminiproblems-problemphyxmini0111-thinlens, def:physics:phyx-mini-0111:phyxminiproblems-problemphyxmini0111-conjugatelocation, def:physics:phyx-mini-0111:phyxminiproblems-problemphyxmini0111-imagenature, def:physics:phyx-mini-0111:phyxminiproblems-problemphyxmini0111-imageorientation, def:physics:phyx-mini-0111:phyxminiproblems-problemphyxmini0111-axialfigurepoint, def:physics:phyx-mini-0111:phyxminiproblems-problemphyxmini0111-raydirectionlabel}
  Physical lenses, figure-derived axial locations, ray-angle readouts, and image classification for the telescope. The eyepiece focal length remains an unknown field; no requested numerical value is built into this structure
\end{definition}
\begin{proof}
  This is the declared model definition of Galilean Telescope Setup.
\end{proof}
\begin{definition}[Matches Galilean Telescope Figure]
  \label{def:physics:phyx-mini-0111:phyxminiproblems-problemphyxmini0111-matchesgalileantelescopefigure}
  \lean{PhyXMiniProblems.ProblemPhyXMini0111.MatchesGalileanTelescopeFigure}
  \uses{def:physics:phyx-mini-0111:phyxminiproblems-problemphyxmini0111-centimetersvalue, def:physics:phyx-mini-0111:phyxminiproblems-problemphyxmini0111-galileantelescopesetup}
  Qualitative and incidence data supplied by the problem and primary figure. The coincident point F₁' = F₂ = I records that I is the objective's would-be image and a virtual object for the diverging eyepiece. The requested eyepiece focal-length readout does not occur here
\end{definition}
\begin{proof}
  This is the declared model definition of Matches Galilean Telescope Figure.
\end{proof}
\begin{definition}[Matches Angular Magnification Readout]
  \label{def:physics:phyx-mini-0111:phyxminiproblems-problemphyxmini0111-matchesangularmagnificationreadout}
  \lean{PhyXMiniProblems.ProblemPhyXMini0111.MatchesAngularMagnificationReadout}
  \uses{def:physics:phyx-mini-0111:phyxminiproblems-problemphyxmini0111-galileantelescopesetup}
  The stated dimensionless angular-magnification readout -6.33
\end{definition}
\begin{proof}
  This is the declared model definition of Matches Angular Magnification Readout.
\end{proof}
\begin{definition}[Has Physical Galilean Parameters]
  \label{def:physics:phyx-mini-0111:phyxminiproblems-problemphyxmini0111-hasphysicalgalileanparameters}
  \lean{PhyXMiniProblems.ProblemPhyXMini0111.HasPhysicalGalileanParameters}
  \uses{def:physics:phyx-mini-0111:phyxminiproblems-problemphyxmini0111-centimetersvalue, def:physics:phyx-mini-0111:phyxminiproblems-problemphyxmini0111-galileantelescopesetup}
  Physical sign and nondegeneracy conditions for the depicted Galilean branch: the objective focal length is positive, the eyepiece focal length is negative, and the incident ray direction has a nonzero small-angle readout
\end{definition}
\begin{proof}
  This is the declared model definition of Has Physical Galilean Parameters.
\end{proof}
\begin{definition}[Satisfies Afocal Galilean Laws]
  \label{def:physics:phyx-mini-0111:phyxminiproblems-problemphyxmini0111-satisfiesafocalgalileanlaws}
  \lean{PhyXMiniProblems.ProblemPhyXMini0111.SatisfiesAfocalGalileanLaws}
  \uses{def:physics:phyx-mini-0111:phyxminiproblems-problemphyxmini0111-lengthreadout, def:physics:phyx-mini-0111:phyxminiproblems-problemphyxmini0111-galileantelescopesetup}
  Paraxial governing laws for an afocal Galilean telescope, stated in every choice of length unit. With signed focal lengths, the convention used by the recorded negative magnification is m = f₁ / f₂, written without division as m * f₂ = f₁. The lens separation is f₁ + f₂, and the angular readouts satisfy u' = m * u
\end{definition}
\begin{proof}
  This is the declared model definition of Satisfies Afocal Galilean Laws.
\end{proof}
\begin{definition}[Has Objective Focal Length Calibration]
  \label{def:physics:phyx-mini-0111:phyxminiproblems-problemphyxmini0111-hasobjectivefocallengthcalibration}
  \lean{PhyXMiniProblems.ProblemPhyXMini0111.HasObjectiveFocalLengthCalibration}
  \uses{def:physics:phyx-mini-0111:phyxminiproblems-problemphyxmini0111-centimetersvalue, def:physics:phyx-mini-0111:phyxminiproblems-problemphyxmini0111-galileantelescopesetup}
  Independent objective calibration needed to recover the recorded answer. The source prose and primary figure do not print this 95 cm value. It is therefore isolated from the actual figure predicate as an explicit additional calibration (for example, a result supplied by omitted surrounding context), rather than being silently attributed to the diagram
\end{definition}
\begin{proof}
  This is the declared model definition of Has Objective Focal Length Calibration.
\end{proof}
\begin{definition}[Answer Choice]
  \label{def:physics:phyx-mini-0111:phyxminiproblems-problemphyxmini0111-answerchoice}
  \lean{PhyXMiniProblems.ProblemPhyXMini0111.AnswerChoice}
  Labels of the four displayed multiple-choice answers
\end{definition}
\begin{proof}
  This is the declared model definition of Answer Choice.
\end{proof}
\begin{definition}[answer Eyepiece Focal Length Centimeters]
  \label{def:physics:phyx-mini-0111:phyxminiproblems-problemphyxmini0111-answereyepiecefocallengthcentimeters}
  \lean{PhyXMiniProblems.ProblemPhyXMini0111.answerEyepieceFocalLengthCentimeters}
  \uses{def:physics:phyx-mini-0111:phyxminiproblems-problemphyxmini0111-answerchoice}
  Eyepiece focal-length readout printed beside each choice, in centimeters
\end{definition}
\begin{proof}
  This is the declared model definition of answer Eyepiece Focal Length Centimeters.
\end{proof}
\begin{definition}[Is Nearest Tenth Readout]
  \label{def:physics:phyx-mini-0111:phyxminiproblems-problemphyxmini0111-isnearesttenthreadout}
  \lean{PhyXMiniProblems.ProblemPhyXMini0111.IsNearestTenthReadout}
  reported is a nearest-tenth decimal readout of exact: it is an integral number of tenths and differs from the exact value by at most half a tenth
\end{definition}
\begin{proof}
  This is the declared model definition of Is Nearest Tenth Readout.
\end{proof}
\begin{lemma}[recorded choice A from objective calibration]
  \label{lem:physics:phyx-mini-0111:phyxminiproblems-problemphyxmini0111-recorded-choice-a-from-objective-calibration}
  \lean{PhyXMiniProblems.ProblemPhyXMini0111.recorded_choice_A_from_objective_calibration}
  \uses{def:physics:phyx-mini-0111:phyxminiproblems-problemphyxmini0111-centimetersvalue, def:physics:phyx-mini-0111:phyxminiproblems-problemphyxmini0111-galileantelescopesetup, def:physics:phyx-mini-0111:phyxminiproblems-problemphyxmini0111-matchesangularmagnificationreadout, def:physics:phyx-mini-0111:phyxminiproblems-problemphyxmini0111-satisfiesafocalgalileanlaws, def:physics:phyx-mini-0111:phyxminiproblems-problemphyxmini0111-hasobjectivefocallengthcalibration, def:physics:phyx-mini-0111:phyxminiproblems-problemphyxmini0111-answereyepiecefocallengthcentimeters, def:physics:phyx-mini-0111:phyxminiproblems-problemphyxmini0111-isnearesttenthreadout}
  If the missing absolute scale is supplied independently as an objective focal length of 95 cm, then the recorded choice A is the nearest-tenth readout. This conditional lemma is not the source-supported main target below
\end{lemma}
\begin{proof}
  The conclusion follows from the governing relations and figure readouts named in the statement of recorded choice A from objective calibration.
\end{proof}
% --- Archon declaration topology end ---
% --- Archon physics formalization source end ---
